% Preamble
\pdfminorversion=4
\documentclass[aspectratio=169]{beamer}
\usepackage{multimedia} % for videos
\usepackage{bm} % for bold math (for tensors)
\usepackage{layout} % put command \layout{} onto an empty frame to see margin lengths
\usepackage{adjustbox}
\usepackage{lipsum}
\usepackage{pdfpages}
\usepackage{hyperref} % for hyperlinks
\usepackage{ulem} % for underline, etc.

% ----- beamer settings -----
% {{{

% define some colors/shades
\definecolor{lightblue}{HTML}{66A3D2}
\definecolor{darkblue}{HTML}{0B61A4}
\definecolor{lightorange}{HTML}{FF8E0B}
\definecolor{darkorange}{HTML}{D45500}
\definecolor{paleorange}{HTML}{FFDEB3}

% set color for hyperlinks
\hypersetup{colorlinks, linkcolor=darkgray, citecolor=darkgray, urlcolor=darkgray}

% set theme
\usetheme{boxes}
\usefonttheme[onlymath]{serif}

% set beamer theme colors, etc.
\setbeamercolor{frametitle}{fg=black}

% customize blocks
\setbeamertemplate{blocks}[rounded][shadow=false]
\setbeamercolor{block title}{fg=white, bg=darkblue}
\setbeamercolor{block body}{fg=black, bg=lightblue}

% customize enumerate and item, including color and spacing
\setbeamercolor*{item}{fg=darkblue}
\setbeamercolor*{enumerate item}{fg=black}
\setbeamercolor*{enumerate subitem}{fg=black}
\setbeamercolor*{enumerate subsubitem}{fg=black}
\setbeamertemplate{itemize subitem}{$-$}
\setlength{\leftmarginii}{10pt}
\newlength\origleftmargini
\setlength\origleftmargini\leftmargini

% customize footer
\setbeamercolor{footline}{fg=gray}
\setbeamertemplate{navigation symbols}{} % get rid of navigation bar

% add page numbers to footer
\makeatletter
\setbeamertemplate{footline}{
    \leavevmode%
    \hbox{%
    \begin{beamercolorbox}[wd=0.5\paperwidth,ht=2.25ex,dp=1ex,left,leftskip=1ex]{}
    \end{beamercolorbox}
    \begin{beamercolorbox}[wd=0.5\paperwidth,ht=2.25ex,dp=1ex,right,rightskip=2ex]{}
      \insertframenumber % page number only
      % \insertframenumber/\inserttotalframenumber % page number/total pages
    \end{beamercolorbox}
    }\vskip0pt
}
\makeatother
% }}}

% ----- tikz ----- (for drawing)
% {{{
\usepackage{tikz}
\usetikzlibrary{shapes,positioning,arrows,backgrounds,fit}
% tikz define invisible: 
% http://tex.stackexchange.com/questions/55806/mindmap-tikzpicture-in-beamer-reveal-step-by-step/55849#55849
\tikzset{
  invisible/.style={opacity=0.2},
  visible on/.style={alt=#1{}{invisible}},
  alt/.code args={<#1>#2#3}{%
    \alt<#1>{\pgfkeysalso{#2}}{\pgfkeysalso{#3}} % \pgfkeysalso doesn't change the path
  },
}

\pgfdeclarelayer{bg}    % declare background layer
\pgfsetlayers{bg,main}  % set the order of the layers (main is the standard layer)

\newcommand{\highlighteq}[2]{ \tikz[baseline]{\node[fill=#1,rounded corners,anchor=base]{$\displaystyle \strut #2$};} }

\tikzset{onslide/.code args={<#1>#2}{%
  \only<#1>{\pgfkeysalso{#2}} % \pgfkeysalso doesn't change the path
}}

% }}}

% --- Notes on options for including movies ---
% {{{

% (1) movie15 package
% \usepackage{movie15}
% \includemovie[text={\includegraphics[width=5cm]{figs/lam_t0.pdf}}]{}{}{figs/lam.avi}
% Comments: 
%   * apparently depreciated (superceded by movie9), but seems to play movies fine
%   * (bad) Puts a pushpin icon over the video
%   * (good) Actually embeds the movie, so don't need to include a separate folder
%   * (bad) Puts black or white bars when playing movie with controls (changes size of
%     the movie to fit the controls into the frame).

% (2) movie9 package
% \usepackage{media9}
% Comments:
%   * Plays using a flash player. Doesn't work on Linux.

% (3) multimedia package
% \usepackage{multimedia}
% \movie[width=5cm, showcontrols]{\includegraphics[width=5cm]{figs/lam_t0.pdf}}{figs/lam.avi}
% Comments:
%   * (good) Seems like the best (but flawed) option for playing a movie in the frame
%   * (bad) Does not embed the movie. Need a separate videos folder.
%   * (bad) Puts black or white bars when playing movie with controls (changes size of
%     the movie to fit the controls into the frame).
%   * (aside) Appears that the backend of the player (same as movie15 pacakge) can be 
%     changed from gstreamer to vlc. Used apt-get install phonon-backend-vlc. Made bars
%     white, not black, which was a marginal improvement.
%     http://tex.stackexchange.com/questions/163089/okular-to-play-embedded-videos-in-beamer

% (4) Hyperlink
% \include{hyperref}
% \href{run:figs/lam.avi}{\includegraphics[width=5cm]{figs/lam_t0.pdf}}
% Comments:
%   * Runs an external viewer (like VLC)
%   * (good) Cross-platform compatible
%   * (bad) Means lots of clicking, and not very easy to move back and forth
%   * (bad) Does not embed the movie. Need a separate videos folder.
%   * (bad) Cannot set up a script to give arguments to external program (e.g. VLC)

% (5) pdfpc
% \href{run:figs/lam.avi?autostart&loop}{\includegraphics[width=5cm]{figs/lam_t0.pdf}}
% Comments:
%   * (neutral) Separate presentation platform. Has nice features.
%   * (bad) Has a slow "pre-rendering" feature.
%   * (bad) All tested videos end up very poor resolution.
%   * (good) Viewer does not re-size videos with ugly controls

% }}}

% ----- tips/tricks & custom definitions ----
% {{{

%% ** block title without a block **
\newcommand{\blocktitle}[1]{
  \begin{tikzpicture}
    \node[ fill, color=darkblue, text=white, rounded corners, inner sep=4pt,
           align=center, minimum width=4cm, minimum height=\baselineskip
    ] 
    {\large #1\par};
  \end{tikzpicture}
}

% ** columns **
%  \begin{columns}[T]
%    \begin{column}{0.45\textwidth}
%      Column 1: Theory
%    \end{column}
%    \begin{column}{0.45\textwidth}
%      Column 2: Experiment
%    \end{column}
%  \end{columns}

% ** make a figure bigger than the textwidth **
%\makebox[\textwidth][c]{\includegraphics[width=1.0\textwidth]{figs/immersion_precipitation_roll.pdf}}

% ** bulk commenting **
%\iffalse
%\fi

% ** Put citation in bottom middle **
%\begin{tikzpicture}[remember picture, overlay]
%  \node[text width=\textwidth, align=center, anchor=south] at (current page.south) {
%    \scriptsize Lee et al. Adv.~Materials. 29, 1076008 (2017)\par
%  };
%\end{tikzpicture}

% }}}

% custom title page
% {{{
\setbeamerfont{author}{size=\normalsize} % font sizes
\newcommand{\location}[1]{\def\insertlocation{#1}}
\newcommand{\titlecolor}{\color{white}}
\newcommand{\bottomcolor}{\color{white}}
% https://tex.stackexchange.com/questions/350511/latex-beamer-create-own-variable
\newcommand\insertpresenter{}
\newcommand\presenter[1]{\renewcommand\insertpresenter{ \uline{#1}} }

\defbeamertemplate*{title page}{customized}[1][]
{
    \begin{tikzpicture}[remember picture, overlay]

      % background photo
      \node [draw=none, align=center] at (current page.center)
      {
        \includegraphics[keepaspectratio, width=1.01\paperwidth]{./figs/cover_art_cropped.pdf}
      };

      % BYU logo
      \node [draw=none, align=center, anchor=north west] 
        at ([xshift=-0.05in, yshift=0.05in] current page.north west)
      {
        \includegraphics[keepaspectratio, width=0.75in]{./figs/BYU_IM22a.pdf}
      };

     % group logo
      \node [draw=none, align=center, anchor=north east] at (current page.north east)
      {
        \includegraphics[keepaspectratio, width=0.6in]{./figs/TreeGroupLogo.pdf}
      };
  
      % title
      \usebeamercolor{title}
      \node [text width=0.80\paperwidth, align=center, draw=none, fill=black, rounded corners,
             fill opacity=0.3, text opacity=1, inner sep=5pt,
             below = 0.3in of current page.center, anchor=center] 
      {
        {
          \huge \titlecolor{} \inserttitle\par
        }
        \vspace{12pt}
        {
          \Large \bottomcolor{}\insertpresenter, \insertauthor \par
        } 
        \vspace{8pt}
        { 
          \bottomcolor{} \insertinstitute \par
        }

      };

      % add date
      \node [draw=none, text width=0.6\paperwidth, anchor=south west] (date)
      at (current page.south west)
      {
        \bottomcolor{}
        \raggedright
        \insertlocation\par
      };

      % add location
      \node [draw=none, above left = 0cm and 0cm of current page.south east, text width=0.6\paperwidth] (date)
      {
        \bottomcolor{}
        \raggedleft
        \insertdate\par
      };

    \end{tikzpicture}
    
}
% }}}

% ----- title/author info -----
\title{Talk Title\\A second line}
\presenter{Speaker Name (underlined)}
% for other authors
\author{Contributing Author1, and Contributing Author2}
\institute{Brigham Young University}
\location{Meeting Name}
\date{Month DD, YYYY}

\begin{document}

\begin{frame}[plain]
  % customize title page using the template in the preamble
  \titlepage
\end{frame}

\begin{frame}[c]{Simple Slide with a Block}
% {{{

  % [c] vertically centers your content
  % [t] will top align the content

  \centering
  \vspace{-0.5\baselineskip}
  \begin{block}{Example Block}
    \begin{itemize} 
    \item This is like Powerpoint
    \item Isn't it a terrible slide?
    \item \href{https://www.overleaf.com/learn/latex/Beamer\_Presentations\%3A\_A\_Tutorial\_for\_Beginners\_(Part\_1)\%E2\%80\%94Getting\_Started}{Here} is a link to a Beamer tutorial by Overleaf. 
It gives some good basics.
  \end{itemize}
  \end{block}

%% }}}
\end{frame}

\begin{frame}{Two Columns with Blocks}
% {{{

  \begin{columns}[T]

  \column{0.45\textwidth}

    \centering
    \vspace{-0.5\baselineskip}
    \begin{block}{Example Block with Items}
      \begin{itemize} 
      \item At least this slide has an image!
      \item Item 2
      \item Item 3
      \item Item 4
    \end{itemize}
    \end{block}

  \column{0.45\textwidth}

    \centering
    \blocktitle{Block Title Only}
    \includegraphics[width=\textwidth]{figs/fig-placeholder.pdf}

    {
      \scriptsize 
      Example citation.\par
      Sosa et al. Macromolecules. 49, 3580 (2016)\par
      You need a par after it to get the text size and spacing to change properly.\par
    }

  \end{columns}

%% }}}
\end{frame}

\begin{frame}{Successively adding material to a slide}
% {{{

  \begin{block}{Using Overlay Specs}
  
  This is called ``Overlay'' in Beamer.
  There are different ways to do this. 
  I give just a few examples.
  See \href{https://www.overleaf.com/learn/latex/Beamer\_Presentations\%3A\_A\_Tutorial\_for\_Beginners\_(Part\_4)\%E2\%80\%94Overlay\_Specifications}{this link} for more.
  Note that the slide number doesn't change.
  \end{block}
  

  \begin{columns}[T]

  \column{0.45\textwidth}

    \centering
    \vspace{-0.5\baselineskip}
    \begin{block}{Example Block with Items}
      \begin{itemize} 
      \item<1-> Item 1
      \item<2,4> Item 2
      \item<3> Item 3
      \item<4-> Item 4
    \end{itemize}
    \end{block}

  \column{0.45\textwidth}

    \centering
    \vspace{-0.5\baselineskip}
  
    \onslide<1->{
    \begin{equation*}
    a^{2} + b^{2} = c^{2}
    \end{equation*}
    }

    \only<2>{
    \begin{equation*}
    e^{-i \pi} = -1
    \end{equation*}
    }

    \onslide<3>{
    \begin{equation*}
    \sin(x)^{2} + \cos(x)^{2} = 1
    \end{equation*}
    }

  \end{columns}

%% }}}
\end{frame}

\begin{frame}<1>[label=outline]{Repeating frames}
% {{{

  \begin{block}{Using Overlay Specs}
  Sometimes you want to repeat a slide.
  A label and the command \texttt{againframe} allows you to do this (even if the frames are not back-to-back).
  \end{block}

  \onslide<1->{
  \begin{equation*}
  a^{2} + b^{2} = c^{2}
  \end{equation*}
  }

  \only<2>{
  \begin{equation*}
  e^{-i \pi} = -1
  \end{equation*}
  }

  \onslide<3>{
  \begin{equation*}
  \sin(x)^{2} + \cos(x)^{2} = 1
  \end{equation*}
  }

%% }}}
\end{frame}

\againframe<2>{outline}

\begin{frame}{Relative Placement Using Tikz}
% {{{

  \begin{block}{Floating Picture with Relative Locations Inside}
  Sometimes you want more control over placement than just columns.
  Tikz is good at this.
  This slides shows Tikz doing relative placement, but the overall TikzPicture is still placed by beamers algorithms to fit with your other content.
  \end{block}

  \begin{tikzpicture}

    \node[text width=4cm, align=center, inner sep=0pt] (A) {
      \includegraphics[width=\textwidth]{figs/fig-placeholder.pdf}
    };

    \node[text width=2cm, align=center, inner sep=0pt, right = 0.3 of A] (B) {
      \includegraphics[width=\textwidth]{figs/TreeGroupLogo.pdf}
    };

    \node[text width=3cm, align=center, inner sep=0pt, right = 0.3 of B] (C) {
      \includegraphics[width=\textwidth]{figs/BYU_IM22a.pdf}
    };

    \node[text width=6.5cm, align=left, inner sep=0pt, 
          below right = 0.1 and 0.1 of A.south, anchor=north west, inner sep=0pt] (D) {
      You can put text in these nodes too\ldots
    };

  \end{tikzpicture}

  \begin{block}{Links to more details about relative positioning}
    \begin{itemize}
    \item \href{https://tex.stackexchange.com/questions/69439/how-can-i-achieve-relative-positioning-in-tikz}{Obligatory Stack Exchange Answer}
    \item \href{https://www.bu.edu/math/files/2013/08/tikzpgfmanual.pdf}{Tikz Documentation}
    \end{itemize}
  \end{block}

% }}}
\end{frame}

\begin{frame}{Absolute Placement Using Tikz}
% {{{

  \begin{tikzpicture}[remember picture, overlay]

    % this stretches the figure to fit, even if it isn't the right ratio.
    % get rid of the adjustbox if you don't want to stretch it.
    \node[inner sep=0pt, align=center, anchor=center, opacity=0.25] (C) at (current page.center) {
      \adjustbox{max width=\paperwidth}{
        \includegraphics[height=\paperheight]{figs/Tree_Group_Photo_2021-04-15.jpg}
      }
    };

    \node[inner sep=0pt, text width=6cm, align=left, anchor=north west] (A) at ([yshift=-1.5cm, xshift=0.5cm] current page.north west) {
      Tikz can also be used to do absolute placement of text and images.
      You can even cover the whole page with an image.
      You can also mix this with relative positioning.
    };

    \node[inner sep=0pt, align=center, anchor=center] (B) at (current page.center) {
      \includegraphics[width=4cm]{figs/fig-placeholder.pdf}
    };

  \end{tikzpicture}

% }}}
\end{frame}

\againframe<3>{outline}

\begin{frame}[plain]
% {{{

  Here is a plain frame with no title and no slide number.

% }}}
\end{frame}

\begin{frame}[t]{Movies}
% {{{

  \centering
  \begin{tikzpicture}

    \node[inner sep=0pt, align=center] (video) {

      \movie[width=7cm, poster, autostart, loop, showcontrols] {
        \includegraphics[width=7cm]{vids/example_video_firstframe.png}}
        {vids/example_video.mp4}

    };

    \node[align=center, text width=0.45\textwidth, right = 0.5 of video] (note)
    {
      \begin{block}{Using the multimedia package}
        We can include movies in the pdf using the multimedia package!
        A movie needs two things stored in the \texttt{vids} directory:
        \begin{itemize}
        \item A figure for the first frame
        \item A movie in \texttt{.avi} or \texttt{.mp4} format
        \end{itemize}
        These movies are not embedded in the pdf, but the relative reference is preserved.
        As such, \texttt{vids} \emph{needs to be moved with the pdf} (this is why I keep them in a separate directory).
      \end{block}
    };

  \end{tikzpicture}

% }}}
\end{frame}

\begin{frame}[t]{Acknowledgements}
%% {{{

  \begin{columns}[T]

    \column{0.52\textwidth}
      \centering
      \begin{tikzpicture}

        \node[text width=\columnwidth, align=center, inner sep=0pt] (A) {
          \includegraphics[width=\textwidth]{figs/Tree_Group_Photo_2021-04-15.jpg}
        }; 

        \node[text width=1cm, align=center, fill=white, rounded corners=0.2cm, 
              opacity=0.5, text opacity=1, 
              above right = 2.5 and 4.2 of A.south west, anchor=south west] (B) {
          \small{}Dakota Banks\par
        };

        \node[text width=1.20cm, align=center, fill=white, rounded corners=0.2cm, 
              draw=lightorange, thick, opacity=0.5, text opacity=1,
              above right = 2.5 and 1.2 of A.south west, anchor=south west] (C) {
          \small{}Rami Alhasan\par
        };

      \end{tikzpicture}

    \column{0.48\textwidth}
      \vspace{-8pt}
      {\small
      \begin{itemize}
        \item {\color{lightorange}Prof. Glenn Fredrickson (UCSB)}
        \item Prof. Hector Ceniceros (UCSB)
        \item {\color{lightorange}Dr. Kris Delaney (UCSB)}
        \item Dr. Tatsuhiro Iwama (Asahi Kasei)
        \item {\color{lightorange}Jan Garcia (UCSB)}
      \end{itemize}
      }

  \end{columns}
  \vspace{12pt}
  \begin{center}
  \begin{tikzpicture}

    \node[inner sep=0pt, align=center ] (A)
    { 
      \includegraphics[width=2.55cm]{./figs/BYU_IM26b.pdf}
    };

    \node[inner sep=0pt, align=center, anchor=center, right=9 pt of A.east] (B)
    { 
      \includegraphics[width=2.05cm]{./figs/ACS-PRF_logo.jpg}
    };
    \node[inner sep=0pt, align=center, text width=3cm, right= 3 pt of B.east] (C) { 
      BYU ORC \\[2pt]
      \includegraphics[width=2.05cm]{./figs/BYU_ORC_chassis.jpg}
    };

  \end{tikzpicture}
  \end{center}



%% }}}
\end{frame}

\end{document}

